% Options for packages loaded elsewhere
\PassOptionsToPackage{unicode}{hyperref}
\PassOptionsToPackage{hyphens}{url}
\PassOptionsToPackage{dvipsnames,svgnames,x11names}{xcolor}
%
\documentclass[
  11pt,
  letterpaper,
]{article}
\usepackage{amsmath,amssymb}
\usepackage{iftex}
\ifPDFTeX
  \usepackage[T1]{fontenc}
  \usepackage[utf8]{inputenc}
  \usepackage{textcomp} % provide euro and other symbols
\else % if luatex or xetex
  \usepackage{unicode-math} % this also loads fontspec
  \defaultfontfeatures{Scale=MatchLowercase}
  \defaultfontfeatures[\rmfamily]{Ligatures=TeX,Scale=1}
\fi
\usepackage{lmodern}
\ifPDFTeX\else
  % xetex/luatex font selection
  \setmainfont[]{DejaVu Serif}
  \setsansfont[]{DejaVu Sans}
  \setmonofont[]{DejaVu Sans Mono}
\fi
% Use upquote if available, for straight quotes in verbatim environments
\IfFileExists{upquote.sty}{\usepackage{upquote}}{}
\IfFileExists{microtype.sty}{% use microtype if available
  \usepackage[]{microtype}
  \UseMicrotypeSet[protrusion]{basicmath} % disable protrusion for tt fonts
}{}
\makeatletter
\@ifundefined{KOMAClassName}{% if non-KOMA class
  \IfFileExists{parskip.sty}{%
    \usepackage{parskip}
  }{% else
    \setlength{\parindent}{0pt}
    \setlength{\parskip}{6pt plus 2pt minus 1pt}}
}{% if KOMA class
  \KOMAoptions{parskip=half}}
\makeatother
\usepackage{xcolor}
\usepackage[margin=1in]{geometry}
\setlength{\emergencystretch}{3em} % prevent overfull lines
\providecommand{\tightlist}{%
  \setlength{\itemsep}{0pt}\setlength{\parskip}{0pt}}
\setcounter{secnumdepth}{-\maxdimen} % remove section numbering
% definitions for citeproc citations
\NewDocumentCommand\citeproctext{}{}
\NewDocumentCommand\citeproc{mm}{%
  \begingroup\def\citeproctext{#2}\cite{#1}\endgroup}
\makeatletter
 % allow citations to break across lines
 \let\@cite@ofmt\@firstofone
 % avoid brackets around text for \cite:
 \def\@biblabel#1{}
 \def\@cite#1#2{{#1\if@tempswa , #2\fi}}
\makeatother
\newlength{\cslhangindent}
\setlength{\cslhangindent}{1.5em}
\newlength{\csllabelwidth}
\setlength{\csllabelwidth}{3em}
\newenvironment{CSLReferences}[2] % #1 hanging-indent, #2 entry-spacing
 {\begin{list}{}{%
  \setlength{\itemindent}{0pt}
  \setlength{\leftmargin}{0pt}
  \setlength{\parsep}{0pt}
  % turn on hanging indent if param 1 is 1
  \ifodd #1
   \setlength{\leftmargin}{\cslhangindent}
   \setlength{\itemindent}{-1\cslhangindent}
  \fi
  % set entry spacing
  \setlength{\itemsep}{#2\baselineskip}}}
 {\end{list}}
\usepackage{calc}
\newcommand{\CSLBlock}[1]{\hfill\break\parbox[t]{\linewidth}{\strut\ignorespaces#1\strut}}
\newcommand{\CSLLeftMargin}[1]{\parbox[t]{\csllabelwidth}{\strut#1\strut}}
\newcommand{\CSLRightInline}[1]{\parbox[t]{\linewidth - \csllabelwidth}{\strut#1\strut}}
\newcommand{\CSLIndent}[1]{\hspace{\cslhangindent}#1}
\usepackage{fontspec}
\usepackage{microtype}
\usepackage{geometry}
\usepackage{xcolor}
\usepackage{graphicx}
\usepackage{eso-pic}
\usepackage{hyperref}
\definecolor{quietgray}{gray}{0.55}
\AtBeginDocument{%
\AddToShipoutPictureFG*{%
  \put(36,22){\makebox[0pt][l]{\scriptsize\color{quietgray}Working Paper · AI-augmented research process}}%
  \put(560,18){\makebox[0pt][r]{\scriptsize\color{quietgray}$\circ\!\!\raisebox{1pt}{\scriptsize ))}$}}%
}}
\setlength{\parindent}{1.2em}
\setlength{\parskip}{0.35em}
\ifLuaTeX
  \usepackage{selnolig}  % disable illegal ligatures
\fi
\usepackage{bookmark}
\IfFileExists{xurl.sty}{\usepackage{xurl}}{} % add URL line breaks if available
\urlstyle{same}
\hypersetup{
  pdftitle={The Shop Makes the Prompt},
  pdfauthor={Watson Hartsoe},
  pdfkeywords={AI art, prompt engineering, cultural
systems, text-to-image, aesthetic competence, human-AI
interaction, generative media},
  colorlinks=true,
  linkcolor={black},
  filecolor={Maroon},
  citecolor={Blue},
  urlcolor={blue},
  pdfcreator={LaTeX via pandoc}}

\title{The Shop Makes the Prompt}
\usepackage{etoolbox}
\makeatletter
\providecommand{\subtitle}[1]{% add subtitle to \maketitle
  \apptocmd{\@title}{\par {\large #1 \par}}{}{}
}
\makeatother
\subtitle{Cultural Competence, Reverse Description, and the Operational
Life of AI Art}
\author{Watson Hartsoe}
\date{2026-08-17}

\begin{document}
\maketitle
\begin{abstract}
Text-to-image prompting is often described as a new language for
expressing artistic intention. That description puts the competence in
the wrong place. Prompting is better understood as a culturally learned
diagnostic practice: practitioners inspect generated artifacts, notice
differences that have become salient through experience, infer what
textual or other intervention might move the system, and test that
intervention against a model whose behavior already contains traces of
training data, platform design, and prior social description. Clifford
Geertz's claim that art and the ``equipment to grasp it'' are made
together becomes unusually literal here. In generative practice,
equipment for grasping an image becomes equipment for correcting the
next one. The prompt is therefore not a transparent statement of
intention and not a stable programming language. It is a provisional
move inside a sociotechnical loop in which description changes causal
position: words that once followed art as criticism can precede art as
conditioning. This paper develops that wager through research on prompt
modifiers, prompt skill, trial-and-error interaction, interface design,
default images, body prompting, Hito Steyerl's ``mean images,'' and
Felipe Rivas San Martín's minority prompt. The consequence is
methodological: the most revealing archive of AI art is not a gallery of
outputs or a collection of final prompts, but a history of noticed
differences, failed generations, corrective moves, hidden priors, and
socially transmitted techniques.
\end{abstract}

\newpage

\section{Introduction: the wrong object is the
prompt}\label{introduction-the-wrong-object-is-the-prompt}

Generative-image culture has made a small text box carry an
extraordinary amount of theory. The box is treated as an interface to
intention, a new artistic medium, a programming language, an
incantation, a compressed brief, and sometimes the location where human
authorship survives. Each description catches something real. Yet all of
them tend to isolate the string that a user submits. They make \emph{the
prompt} look like the stable unit of creative action.

The evidence in this field points elsewhere. Research on text-to-image
practice finds that effective prompting is learned through
experimentation; prompt modifiers circulate as practical techniques;
inexperienced users can recognize prompt quality while still lacking the
style-specific vocabulary needed for effective refinement; and
open-ended text interaction often produces brute-force trial and error
when results are poor (Oppenlaender 2022; Oppenlaender, Linder, and
Silvennoinen 2023; Liu and Chilton 2021). Interface design also changes
how much users continue to explore: shortcuts for producing variants are
associated with reduced exploration of novel concepts and with less
detailed prompting (Torricelli et al. 2023). These findings make the
final string a misleadingly thin artifact. They point to a trajectory in
which a person learns what to notice, which difference matters, and what
kind of intervention might change it.

This paper makes a stronger claim. \textbf{Prompt expertise is not
fundamentally mastery of a language. It is the acquisition of culturally
and technically situated equipment for noticing correctable
differences.} The user does not merely translate an intention into
words. The user encounters an output, recognizes a mismatch through
learned aesthetic distinctions, forms a causal hypothesis about a partly
opaque generator, and tries another move. What appears as linguistic
skill is inseparable from perceptual training, platform affordances,
model-specific regularities, community lore, and the sedimented social
descriptions contained in training data.

Geertz provides the hinge. In \emph{Art as a Cultural System}, he argues
that aesthetic capacities do not stand outside the worlds in which
artworks are made; art and the capacities required to grasp it are
formed together (Geertz 1976). The generative case makes that
proposition operational. The equipment for grasping an AI image is often
immediately reused as equipment for changing the next image. Seeing
becomes a control problem. Description migrates from commentary into
conditioning.

That migration is not simply a triumph of linguistic agency. It also
exposes asymmetry. The generator can fall back toward recurring default
images when textual conditioning fails to discriminate strongly (Simonen
et al. 2025). Training data can render social averages as if they were
machine-native aesthetics (Steyerl 2023). Artists can be forced to spend
prompt effort counteracting class, race, gender, or archival biases that
originated upstream, as Rivas San Martín's ``minority prompt'' makes
explicit (Rivas San Martín 2025). The competence acquired by
practitioners is therefore partly competence in the system's failures.

The argument proceeds in five steps. First, prompting is reconstructed
as diagnostic practice rather than string production. Second, Geertz's
``equipment to grasp'' is extended cautiously into an equipment to
correct. Third, prompt language is shown to be a reverse description of
an inherited image-text archive rather than a direct description of the
pictured world. Fourth, the paper separates user skill from the cultural
and infrastructural priors that skill must work against. Finally, it
shows why description itself changes causal position in generative
systems, producing a feedback loop in which criticism, community
discourse, and interface choices can become production infrastructure.

\section{1. Prompting as diagnostic
practice}\label{prompting-as-diagnostic-practice}

Oppenlaender's ethnography of early text-to-image practice identifies
six classes of prompt modifier and describes prompt engineering as an
iterative, experimental activity (Oppenlaender 2022). The importance of
that finding is not the taxonomy alone. A modifier is learned in
relation to observed output. The practitioner tries a phrase, inspects
what happened, repeats or alters the experiment, and circulates what
seems to work. Prompt terms therefore acquire an \emph{operational
reputation} in addition to ordinary semantic meaning.

This distinction matters because operational reputation can be wrong.
Oppenlaender explicitly discusses idiosyncratic practitioner choices and
the possibility of folk theories about causation (Oppenlaender 2022).
Stochastic generation makes a dangerous epistemic environment: a salient
successful output can follow a recent prompt change without having been
robustly caused by it. A community can stabilize the technique socially
even when the causal theory is weak. Thus the same process that produces
expertise can also produce superstition. The correct research question
is not whether prompt lore exists, but which parts are computationally
effective, under what model versions, and which parts mainly organize
community identity or expectation.

The experimental literature on prompt skill reinforces the process view.
Oppenlaender, Linder, and Silvennoinen found that untrained participants
could evaluate prompt quality and produce descriptive prompts, yet
lacked the style-specific vocabulary necessary for more effective
prompting (Oppenlaender, Linder, and Silvennoinen 2023). This result is
often summarized as evidence that prompt engineering is a learnable
skill. More interestingly, it leaves open \emph{what is actually
learned}. A vocabulary can be memorized. A skilled practice, however,
requires knowing when a distinction is relevant and when a term is worth
trying.

Liu and Chilton locate the problem at the interface level. They describe
open-ended text interaction as double-edged: it offers enormous
expressive freedom but can force users into brute-force trial and error
when a generation fails (Liu and Chilton 2021). The result complicates
claims that natural language removes formalization. The formal work may
simply occur later. A vague initial description produces an artifact;
the artifact makes a missing constraint visible; the user adds that
constraint; the next artifact reveals another. Specification accumulates
temporally through failure.

This suggests a different unit of analysis:

\begin{quote}
output $\rightarrow$ noticed difference $\rightarrow$ causal hypothesis $\rightarrow$ intervention $\rightarrow$ new output.
\end{quote}

The prompt string is only one state in this loop. What changes the
research problem is the \emph{noticed difference}. Two people can
receive the same image and possess different actionable worlds because
one sees a lighting mismatch, a compositional cliché, a model-specific
hand failure, an unwanted social stereotype, or a telltale default that
the other does not recognize. The deepest evidence of expertise may
therefore appear one moment before the next prompt is written.

\section{2. From equipment to grasp to equipment to
correct}\label{from-equipment-to-grasp-to-equipment-to-correct}

Geertz's argument about art resists the idea that aesthetic perception
is a universal capacity applied to autonomous objects. The ability to
respond to art is itself cultivated in forms of life, and his memorable
formulation that art and the equipment to grasp it are made together
gives us a better way to define the relevant cultural unit (Geertz
1976). The question is not necessarily whether ``AI art'' is one
culture. It is where shared sensitivities are reproduced.

For generative practice, those sensitivities are unusually consequential
because they can immediately become operations. A practitioner learns to
see a difference, names it, and feeds that name back into generation.
The cultural formation of perception becomes part of a control loop. We
can therefore split expertise into at least three coupled capacities:

\begin{enumerate}
\def\labelenumi{\arabic{enumi}.}
\tightlist
\item
  \textbf{discrimination} - noticing a difference in an output;
\item
  \textbf{diagnosis} - forming a hypothesis about what produced the
  difference;
\item
  \textbf{intervention} - choosing a prompt, parameter, body movement,
  or interface action intended to alter the next state.
\end{enumerate}

This is stronger than saying experts know more prompt terms. A novice
can copy a phrase without knowing when it applies. Conversely, an expert
may recognize a problem without possessing a reliable intervention
because the model has changed. The mapping from difference to
intervention is versioned and local.

The distinction also explains why prompt secrecy matters. Early prompt
communities sometimes withheld prompts for commercial reasons
(Oppenlaender 2022). If outputs circulate publicly while procedures
remain private, the cultural ``shop'' can split. A public sphere may
teach spectators what to admire while a restricted operative sphere
controls how those effects are reproduced. Aesthetic competence becomes
stratified: perception may be broadly distributed while generative
competence remains scarce.

This is one reason to resist the metaphor of a universal prompt
language. Languages imply some durable relation between expression and
meaning. Prompt practice is closer to a field of situated
correspondences among words, models, interfaces, versions, and
communities. The same phrase can retain ordinary-language meaning while
losing its generative effect after a model update. Conversely, a
seemingly strange phrase can remain operationally useful because it
indexes learned correlations that have little to do with literal
description.

\section{3. Reverse description: prompting the archive that described
the
world}\label{reverse-description-prompting-the-archive-that-described-the-world}

The most consequential complication comes from the training relation
between language and images. Oppenlaender notes that effective
practitioners may have to imagine how other people on the Web would have
described or reacted to an image (Oppenlaender 2022). This changes the
direction of description. The user is not always naming the desired
visual world directly. The user may be estimating the language that
historically surrounded visually similar material in the training
ecology.

Prompting can therefore contain an inverse problem:

\begin{quote}
desired visual tendency $V$ $\rightarrow$ hypothesize historical wording $L$ $\rightarrow$ submit $L$ $\rightarrow$ model maps toward $V$.
\end{quote}

The competence here is partly archival, though the archive is
inaccessible and statistical rather than a conventional catalog. Terms
such as artist names, media labels, genre descriptors, camera
vocabulary, platform names, and evaluative phrases can work because of
how images and texts were co-produced online. The user learns the
aftereffects of those associations by experiment.

Steyerl's ``mean images'' supplies a cultural theory of what sits
upstream of this practice. Generative images, she argues, are
statistical renderings of socially produced data; the apparent machine
image can be a ``social filter'' rendering correlated averages and
latent social patterns (Steyerl 2023). This breaks the usual opposition
between machine style and user taste. There is at least a third term:
the historical distribution of images, descriptions, values, and
classifications already sedimented into the training process.

This matters for Geertz. A naïve application of cultural context to
prompting says that adding ``Edo-period,'' ``queer archive,'' or another
culturally specific label injects cultural meaning. That is not enough.
A cultural name can increase visual recognizability while decreasing
contextual thickness. It may index a statistical cluster of recognizable
markers rather than situated knowledge. The generated surface can look
more specific precisely because it has collapsed local distinctions into
a portable stereotype.

The research problem is therefore not ``does the prompt contain cultural
context?'' It is: \textbf{which cultural distinctions are available to
the model, which are available to the practitioner, and which are
available only to participants in the living practice being invoked?}
These three sets need not coincide.

\section{4. The user corrects what the user did not
choose}\label{the-user-corrects-what-the-user-did-not-choose}

If prompt skill develops partly as competence in a model's inherited
associations, then expertise has a political asymmetry. Some of the work
users learn to perform exists because the system's prior is not neutral.

Rivas San Martín's \emph{Inexistent Archive} provides a precise case.
The project imagines fictional historical photographs of queer,
non-binary, and working-class people in Latin America. Rivas describes
the need to counteract class and race biases in model training data and
develops the concept of the ``minority prompt'' from this obstacle
(Rivas San Martín 2025). The prompt becomes more than description. It is
a local corrective operation against upstream representational
conditions.

That correction should not be confused with structural repair. If a user
adds counter-conditioning to obtain one desired result while the model
and dataset remain unchanged, the user has altered the local generation,
not necessarily the system that made the correction necessary. The labor
is downstream; authority over the causal substrate is upstream. Prompt
agency and prompt burden are therefore compatible.

The same project makes another inversion visible. Rivas retains bodily
errors as an ethical-political limit that marks the images' fabricated
origin and prevents the speculative archive from covering over the
violence that prevented those records from existing (Rivas San Martín
2025). Technical failure becomes provenance. A model improvement that
eliminates malformed bodies can consequently remove an ethical
disclosure device. ``Better'' generation is not monotonically better
art.

Simonen and colleagues' work on default images gives a different form of
upstream pressure. Their study shows that text-to-image systems can
produce visually similar outputs across unrelated or unknown prompts and
analyzes this phenomenon across more than 750,000 Midjourney images
(Simonen et al. 2025). A default is therefore not only a UI setting
selected before interaction. It can be an observable fallback behavior
exposed when language fails to provide enough discriminating guidance.
Deliberately poor prompts can act as probes into the generator's
attractors.

Together, the minority prompt and default image reveal two opposite
encounters with the prior. In one, the user adds language to push
against an unwanted tendency. In the other, language stops steering and
the tendency becomes visible. Both suggest that a complete archive of
prompting must preserve more than successful final strings. It needs
baseline behavior, failures, counter-prompts, version information, and
the evidence by which a user diagnosed the problem.

\section{5. Interfaces choose the cheap next
move}\label{interfaces-choose-the-cheap-next-move}

The model is not the only place where possibility is shaped. Torricelli
and colleagues' analysis of more than 145,000 prompts across two
generative platforms finds that interfaces offering shortcuts for image
variants and diverting attention from prompt editing are associated with
reduced exploration of novel concepts and less detail in prompts
(Torricelli et al. 2023). The result supplies a concrete mechanism for
platform aesthetics that does not require assuming an opaque ranking
algorithm.

A platform can change creative trajectories simply by changing the cost
of the next action. If ``make variants'' is one tap while re-description
requires more effort, local exploitation becomes cheaper than conceptual
movement. The interface does not need to impose a style explicitly. It
can alter transition probabilities through affordances.

This distinction matters because claims about ``AI style'' frequently
collapse multiple causal layers: model architecture, training data,
prompt population, interface actions, ranking, and imitation. The
evidence currently supports some mechanisms more strongly than others.
Interface-mediated exploration has observational support (Torricelli et
al. 2023). The stronger claim that a particular platform ranking system
causes a specific visual style requires its own receipts. A
cultural-systems account should not smooth these layers into a
harmonious network; it should force them to make different predictions.

The practical consequence is methodological. To study generative
aesthetics, researchers should hold layers constant wherever possible.
Keep the model fixed and change interface actions. Keep prompts fixed
and change models. Keep architecture fixed and change training
distributions where reproducible models permit it. Preserve the
mismatches. ``Culture'' becomes analytically useful when it helps locate
mechanisms rather than when it simply names everything surrounding the
image.

\section{6. Description changes causal
position}\label{description-changes-causal-position}

Susan Sontag's call for a descriptive rather than over-interpretive
criticism creates an unexpected bridge to generative practice. She
wanted criticism to show how an artwork is what it is, resisting the
reduction of sensuous form to hidden content (Sontag 1966). In ordinary
criticism, the arrow runs from form to words. A critic observes
contrast, texture, framing, rhythm, scale, or tone and develops language
adequate to what is already there.

In a text-conditioned generator, much of the same vocabulary can move
upstream:

\begin{quote}
FORM $\rightarrow$ WORD \hspace{1cm} becomes \hspace{1cm} WORD $\rightarrow$ CONDITIONING $\rightarrow$ FORM.
\end{quote}

The word does not deterministically compile into the visual feature, but
it becomes causally operative. A descriptive vocabulary becomes a
control surface. This is a deeper transformation than simply ``using
words to make pictures.'' Cultural vocabularies for noticing art can
become ingredients in the statistical reproduction of future art.

Yet even here, text should not be mistaken for the essence of prompting.
Oppenlaender and colleagues' body-prompting installation demonstrates
generative conditioning through embodied input in a public art setting
(Oppenlaender et al. 2024). The category ``prompt'' survives after
textual language disappears. What unifies text prompting and body
prompting is not syntax but intervention into the system's next state.

That observation disciplines the stronger metaphor that ``the body is a
language.'' An input modality is not automatically a language. To
establish an embodied syntax we would need evidence of stable units,
recombination, learned correspondences, malformed combinations, or
systematic semantics. The source establishes control, not grammar. The
larger lesson is useful: prompting should be defined by its position in
a generative transition before being defined by its representational
medium.

\section{7. The cultural loop without model
learning}\label{the-cultural-loop-without-model-learning}

A final consequence follows from the circulation of prompt knowledge.
Community discourse does not have to change model weights in order to
change what a fixed model produces. Outputs are posted; prompts are
disclosed or reverse-engineered; aesthetic judgments attach to them;
terms circulate; another user inserts those terms into a new prompt; and
a new output enters circulation (Oppenlaender 2022). The cultural loop
can operate around a technically fixed generator.

This gives a sharper account of ``meaning-in-use'' for generative
systems. Use can mean interpretation, but interpretation can become
future input through social transmission. Reception becomes production
infrastructure without any adaptive model update. What changes is the
human input distribution.

The same loop also explains why technically false prompt theories can
matter. A modifier may have weak causal effect yet strong social effect
if it signals expertise, taste, affiliation, or adherence to a community
recipe. Conversely, a technically powerful modifier may remain
culturally invisible. Prompt expressions can operate on two state spaces
at once: machine conditioning and social organization. The two effects
should be measured separately.

\section{Discussion: archive the difference, not just the
string}\label{discussion-archive-the-difference-not-just-the-string}

The paper's wager can now be stated compactly. \textbf{The most
consequential unit of prompt culture is not the prompt term but the
learned mapping from a noticed difference to a possible intervention.}
This mapping is culturally acquired, technically contingent, and
politically uneven.

That claim changes what should be archived. A final prompt suppresses
the reason each phrase entered. A final image suppresses rejected
alternatives. A prompt guide suppresses failed causal hypotheses. A
screenshot suppresses model version, interface action costs, and
defaults. If scholars want to understand generative art as a cultural
system, they need trajectories that preserve at least:

\begin{itemize}
\tightlist
\item
  the desired state as understood at that moment;
\item
  the generated candidate;
\item
  the difference the practitioner noticed;
\item
  the explanation the practitioner entertained;
\item
  the intervention chosen;
\item
  the model and interface state;
\item
  whether the intervention worked across repetitions;
\item
  how the technique was learned, shared, withheld, or contested.
\end{itemize}

This is not an argument for total provenance. Total provenance is
impossible and can become its own fetish. It is an argument that the
culturally interesting object is often the transition by which a
difference becomes actionable.

The proposal also yields direct empirical tests. Expert and novice
participants can be shown identical flawed generations and asked to
annotate every discrepancy before editing the prompt. Their noticed
differences can be compared before their vocabularies are compared.
Prompt folklore can be ablated across seeds and model versions. Matched
interfaces can vary only the cost of variant generation versus
re-description. Cultural labels can be evaluated by practitioners from
the invoked traditions rather than by generic recognizability.
Default-image probes can test what surfaces when language ceases to
discriminate.

\section{Limitations and unresolved
territory}\label{limitations-and-unresolved-territory}

Several boundaries matter. First, the field joins sources produced
across different generations of text-to-image systems. Practices
documented around VQGAN-CLIP-era tools cannot automatically be
generalized to contemporary diffusion and multimodal systems. The
instability is part of the object: operational prompt knowledge can
expire.

Second, the proposed ``equipment to correct'' is an extension of Geertz,
not his terminology. Geertz supplies an account of culturally formed
aesthetic competence; this paper asks what happens when that competence
enters an iterative generative loop. Historical influence is not
claimed.

Third, the argument does not show that all prompt skill is perceptual.
Some expertise may be lexical, technical, strategic, social, or
domain-specific. The stronger claim is that vocabularies alone cannot
explain effective correction without attention to the distinctions users
learn to notice.

Fourth, corpus effects, model architecture, interface design, ranking,
and current user taste remain difficult to separate causally. Steyerl's
``mean image'' is a critical description of socially sedimented
statistical rendering, not a quantitative decomposition of those causes
(Steyerl 2023). The field needs controlled interventions rather than
broader synthesis.

Finally, not every meaningful use of generative art is aimed at
correction. Body prompting, chance operations, intentional misuse, and
practices that cultivate surprise can reject the goal of converging on a
pre-specified image. The diagnostic loop is strongest where
practitioners are steering toward or away from recognizable conditions.
Its boundary is exactly where ``failure'' ceases to be a defect and
becomes the event the work was seeking.

\section{Conclusion: the shop makes the
prompt}\label{conclusion-the-shop-makes-the-prompt}

A prompt appears to begin with words. The research reviewed here
suggests that it begins earlier, in a learned capacity to see what those
words might need to change.

Geertz's shop is therefore not a metaphorical decoration for AI art. It
identifies a concrete research problem. The community makes images, but
it also makes the sensitivities by which images are judged; those
sensitivities make failures legible; failures motivate interventions;
interventions circulate as prompt lore; and that lore changes what the
fixed generator is asked to produce. Meanwhile the generator carries its
own inherited social distributions, defaults, and biases, forcing users
to learn not only how to describe worlds but how to negotiate the
machine's sedimented descriptions of worlds.

The prompt is not the cultural system. It is one move through it. The
better object is the loop in which perception becomes correction and
description becomes operation.

\section{Appendix A - Assembly
Instrument}\label{appendix-a---assembly-instrument}

The exact assembly instrument for this working paper is preserved
verbatim as \texttt{SLIPCASE\_FINAL\_PROMPT.txt} and
\texttt{\_PROMPTS/assembly\_prompt\_\_v15.55-AM.txt}. The complete
prompt is not reproduced in the paper body because the package and
standalone \texttt{index.html} preserve it directly.

\section{Appendix B - Making History}\label{appendix-b---making-history}

This paper was assembled from forty compiled zettels, three
user-provided research PDFs, the supplied recursive forage instrument,
and public source records and PDFs registered in the accompanying
resource ledger. The zettel payloads in this checkpoint were
reconstructed from visible conversation output rather than copied from a
machine-readable full-chat export; their hashes therefore identify the
compiled checkpoint payloads, not an unavailable upstream transcript.
Full control, retrieval, preservation, and uncertainty notes are
recorded in \texttt{000\_\_MAKING\_HISTORY.txt}.

\section{Appendix C - Replication
Path}\label{appendix-c---replication-path}

Checkpoint: \texttt{SLIPCASE-20260817-AIACS-01}. Begin with
\texttt{000\_\_RETURN\_PATH.txt} and \texttt{000\_\_REBUILD.txt}. The
claim-to-evidence chain is in
\texttt{the-shop-makes-the-prompt\_\_SOURCE\_MAP.txt}; machine state is
under \texttt{\_SLIPCASE/}; exact zettel payloads are root \texttt{.txt}
cards with byte-identical \texttt{\_MD/} mirrors. Recompute hashes
before any merge and preserve unresolved addresses as ghosts rather than
fuzzy-matching them.

\section*{References}\label{references}
\addcontentsline{toc}{section}{References}

\phantomsection\label{refs}
\begin{CSLReferences}{1}{0}
\bibitem[\citeproctext]{ref-Geertz1976ArtCulturalSystem}
Geertz, Clifford. 1976. {``Art as a Cultural System.''} \emph{MLN} 91
(6): 1473--99. \url{https://doi.org/10.2307/2907147}.

\bibitem[\citeproctext]{ref-LiuChilton2021PromptEngineering}
Liu, Vivian, and Lydia B. Chilton. 2021. {``Design Guidelines for Prompt
Engineering Text-to-Image Generative Models.''}
\url{https://doi.org/10.48550/arXiv.2109.06977}.

\bibitem[\citeproctext]{ref-Oppenlaender2022PromptModifiers}
Oppenlaender, Jonas. 2022. {``A Taxonomy of Prompt Modifiers for
Text-to-Image Generation.''}
\url{https://doi.org/10.48550/arXiv.2204.13988}.

\bibitem[\citeproctext]{ref-OppenlaenderEtAl2024BodyPrompting}
Oppenlaender, Jonas, Hannah Johnston, Johanna Silvennoinen, and Helena
Barranha. 2024. {``Artworks Reimagined: Exploring Human-AI Co-Creation
Through Body Prompting.''}
\url{https://doi.org/10.48550/arXiv.2408.05476}.

\bibitem[\citeproctext]{ref-OppenlaenderLinderSilvennoinen2023PromptingAIArt}
Oppenlaender, Jonas, Rhema Linder, and Johanna Silvennoinen. 2023.
{``Prompting AI Art: An Investigation into the Creative Skill of Prompt
Engineering.''} \url{https://doi.org/10.48550/arXiv.2303.13534}.

\bibitem[\citeproctext]{ref-Rivas2025InexistentArchive}
Rivas San Martín, Felipe. 2025. {``Inexistent Archive.''} \emph{ReCIA -
Journal of the Arts Research Centre}, no. 1: 153--71.
\url{https://doi.org/10.21134/a3y8jg10}.

\bibitem[\citeproctext]{ref-SimonenEtAl2025DefaultImages}
Simonen, Hannu, Atte Kiviniemi, Hannah Johnston, Helena Barranha, and
Jonas Oppenlaender. 2025. {``An Exploration of Default Images in
Text-to-Image Generation.''}
\url{https://doi.org/10.48550/arXiv.2505.09166}.

\bibitem[\citeproctext]{ref-Sontag1966AgainstInterpretation}
Sontag, Susan. 1966. \emph{Against Interpretation and Other Essays}.
Farrar, Straus; Giroux.

\bibitem[\citeproctext]{ref-Steyerl2023MeanImages}
Steyerl, Hito. 2023. {``Mean Images.''} \emph{New Left Review}, no.
140/141: 82--97. \url{https://doi.org/10.64590/uhm}.

\bibitem[\citeproctext]{ref-TorricelliEtAl2023Interface}
Torricelli, Maddalena, Mauro Martino, Andrea Baronchelli, and Luca Maria
Aiello. 2023. {``The Role of Interface Design on Prompt-Mediated
Creativity in Generative AI.''}
\url{https://doi.org/10.48550/arXiv.2312.00233}.

\end{CSLReferences}

\end{document}
